\section*{الفصل \ref{ch:b3:milky-way-galaxies} --- درب التبانة والمجرّات}

\begin{solution}{exo:b3:milky-way-galaxies:1}
(أ)~$2\pi R_0 = 2\pi\times\qty{2.523e17}{km} = \qty{1.585e18}{km}$؛ عند
\qty{229}{km.s^{-1}} يستغرق هذا \qty{6.92e15}{s}، أي
\qty{219}{Myr}. (ب)~$4540/219 = 20.7$: نحو عشرين دورة.
\end{solution}

\begin{solution}{exo:b3:milky-way-galaxies:2}
(أ)~$\qty{1}{au}/\qty{1}{yr} = \num{1.496e8}/\num{3.156e7} =
\qty{4.74}{km.s^{-1}}$. (ب)~$v = 4.74\times6.411\times8.178 =
\qty{248.5}{km.s^{-1}}$، ضمن ما تُقدّمه المراجعة
\qty[separate-uncertainty=true]{248 \pm 3}{km.s^{-1}}. (ج)~$248.5 - 11 = \qty{237.5}{km.s^{-1}}$،
أعلى بنحو \qty{8}{km.s^{-1}} من \qty{229}{km.s^{-1}} لدى العمالقة الحمراء:
يرتكز الرقمان على قيمتين مختلفتين للمسافة $R_0$ (8.178 و\qty{8.122}{kpc})،
\omterm{def:b3:milky-way-galaxies:lsr}{ومعيار السكون المحلي} نفسه قد ينساب بمقدار يصل إلى نحو \qty{15}{km.s^{-1}} بالنسبة
إلى مدار دائري حقيقي.
\end{solution}

\begin{solution}{exo:b3:milky-way-galaxies:3}
(أ)~$n = 10(1 - 0.6) = 4$: E4. (ب)~$b/a = 1 - 7/10 = 0.3$. (ج)~يعتمد التصنيف على
نسبة المحورين الظاهرية؛ فالمجسّم الإهليلجي المفلطح المرئي على طول محوره القصير
له مخطّط دائري ويُصنَّف E0.
\end{solution}

\begin{solution}{exo:b3:milky-way-galaxies:4}
\mbox{\textdir TLT $d = 10^{\mu/5 + 1}$~pc}: سحابة ماجلان الكبرى \qty{49.6}{kpc}، أندروميدا
\qty{783}{kpc}، \babelsublr{NGC~4258} \qty{7.58}{Mpc}. لنجمٍ في السحابة
اختلاف منظر قدره $1/\qty{49.6}{kpc} = \qty{0.020}{mas}$، وهو
مقدار عدم اليقين في أفضل كتالوج للنجوم الساطعة: فاختلاف المنظر الواحد لا يعطي
أيّ مسافة على الإطلاق، ويُحدَّد موضع السحابة بوسائل أخرى (الثنائيات الكاسفة)،
ومنها يبدأ \omterm{def:b3:milky-way-galaxies:ladder}{سلّم المسافات} القيفاوي.
\end{solution}

\begin{solution}{exo:b3:milky-way-galaxies:5}
(أ)~$\omega_0 = A - B = \qty{27.2}{km.s^{-1}.kpc^{-1}}$ و$V_0 =
27.2\times8.178 = \qty{222}{km.s^{-1}}$. (ب)~$V'_0 = -(A + B) =
\qty{-3.4}{km.s^{-1}.kpc^{-1}}$: منحنى يهبط ببطء. (ج)~كلتاهما
أدنى من قيم العمالقة الحمراء (229 و\qty{-1.7}{km.s^{-1}} لكلّ
كيلوفرسخ فلكي)؛ وجد مؤلفو ذلك المنحنى النتيجة نفسها، فقيمة $V_0$ أدنى بنحو
انحرافين معياريين بالنسبة إلى $R_0$ وثابتَي أورت هذين. صيغُ أورت
تفترض مدارات دائرية في قرص متماثل محوريًّا؛ والثابتان $C$ و$K$ غير الصفريّين
يدلّان على أنّ حقل السرعة المحلي يتأثّر بالعصابة والأذرع \omterm{def:b3:milky-way-galaxies:types}{الحلزونية}، وأنّ
الثابتان المحليّان $A$ و$B$ يحملان تلك الاضطرابات.
\end{solution}

\begin{solution}{exo:b3:milky-way-galaxies:6}
(أ)~$\rho = V^2/(4\pi GR^2) = 229^2/(4\pi\times\num{4.301e-6}\times
8.178^2) = \num{1.45e7}$ كتلةً شمسية لكلّ كيلوفرسخ فلكي مكعّب، أو
\num{0.0145} لكلّ فرسخ فلكي مكعّب. (ب)~$M_\odot c^2 = \num{1.9884e30}\times
(\num{2.998e8})^2 = \qty{1.787e47}{J} = \qty{1.115e57}{GeV}$
و$\qty{1}{pc^3} = (\qty{3.0857e18}{cm})^3 = \qty{2.938e55}{cm^3}$: والنسبة هي
\qty{37.96}{GeV.cm^{-3}}. ثمّ $0.30/37.96 = \num{0.0079}$ كتلة شمسية
لكلّ فرسخ فلكي مكعّب. (ج)~تُحصي كثافةُ المنحنى المسطّح كلَّ الكتلة التي تُمسك
بالمدارات، بما فيها النجوم والغاز، كأنّها موزَّعة في كرة؛
والمادة المظلمة ليست إلّا جزءًا منها، نحو نصفها هنا.
\end{solution}

\begin{solution}{exo:b3:milky-way-galaxies:7}
(أ)~$W_I(\mathrm{LMC}) = -3.314\times1.4771 + 15.888 = \qty{10.993}{mag}$،
أي $M_W = 10.993 - 18.477 = \qty{-7.484}{mag}$. (ب)~$\mu = 25.00 + 7.484 =
32.484$ و$d = 10^{32.484/5 + 1}~\mathrm{pc} = \qty{31.4}{Mpc}$.
(ج)~$\delta d/d = (\ln 10/5)\,\delta\mu = 0.4605\times0.077 =
\qty{3.5}{\percent}$ لنجمٍ واحد، و$3.5/\sqrt{50} = \qty{0.5}{\percent}$
لخمسين نجمًا، إذا كانت أخطاؤها مستقلّة (إذ يُضاف خطأ مسافة السحابة نفسها
إلى الجميع بالقدر نفسه).
\end{solution}

\begin{solution}{exo:b3:milky-way-galaxies:8}
$\mu = I - A_I - M_I = 26.00 - 0.05 + 4.05 = 30.00$، أي $d =
10^{7}~\mathrm{pc} = \qty{10.0}{Mpc}$. خطأ قدره \qty{0.04}{mag} يساوي
$0.4605\times0.04 = \qty{1.8}{\percent}$ في المسافة.
\end{solution}

\begin{solution}{exo:b3:milky-way-galaxies:9}
(أ)~$d = cz/H_0$ يعطي $5\log_{10}d = 5\log_{10}(cz) - 5\log_{10}H_0$،
أي $m_B - M_B - 25 = 5\log_{10}(cz) - 5\log_{10}H_0$؛ ومع
$\log_{10}(cz) = a_B + 0.2\,m_B$، و$5\log_{10}H_0 = 5a_B + M_B + 25$، وهي
النتيجة المطلوبة. (ب)~$\log_{10}H_0 = -3.8506 + 0.714158 + 5 = 1.8636$، أي
$H_0 = \qty{73.04}{km.s^{-1}.Mpc^{-1}}$. (ج)~$M_B = 5(\log_{10}67.4 -
5.714158) = \qty{-19.427}{mag}$: المستعرات الأعظم أكثر سطوعًا بمقدار \qty{0.174}{mag}،
وبذلك تكون كلّ المسافات أكبر بعامل $10^{0.2\times0.174} = 1.084$، أي
نسبة $73.04/67.4$. ويمثّل التوتّر خلافًا بنحو ثمانية في المئة في
مقياس المسافات.
\end{solution}

\begin{solution}{exo:b3:milky-way-galaxies:10}
(أ)~نجمٌ عند الحافة، في مدار دائري، يحقّق $V^2/R \approx GM/R^2$
(تقريب كروي). (ب)~$R = \sqrt{M/(\pi\Sigma)}$، أي $V^2 =
GM/R = G\sqrt{\pi\Sigma}\,M^{1/2}$ و$M = V^4/(\pi\Sigma G^2) \propto
V^4$؛ مع نسبة مشتركة بين الكتلة والضوء، $L \propto V^4$ أيضًا.
(ج)~$(100/229)^4 = \num{0.036}$: المجرّة الأسرع تزيد كتلتها
بمقدار 27 مرّة. (د)~عند تساوي الإشعاعية يتناسب الفيض مع $1/d^2$، فالمجرّة الأخفت
تبعد ضعف المسافة.
\end{solution}

\begin{solution}{exo:b3:milky-way-galaxies:11}
(أ)~$\qty{1}{km.s^{-1}} = \qty{1.0227}{kpc.Gyr^{-1}}$، أي
\qty{229}{km.s^{-1}} تساوي \qty{234.2}{kpc.Gyr^{-1}}. (ب)~$2\pi R/V$:
\qty{107}{Myr} عند \qty{4}{kpc}، و\qty{322}{Myr} عند \qty{12}{kpc}.
(ج)~$\Delta\phi = Vt\,(1/4 - 1/12) = 234.2\times0.1667 = \qty{39.0}{rad}$،
أي نحو 6.2 دورة: سيكون الذراع لولبًا من ست لفّات محكمة بين 4
و\qty{12}{kpc}، أضيق بكثير من الأذرع المرصودة، وأضيق أكثر عند كلّ
دورة أخرى. (د)~في موجة الكثافة يدور النمط بسرعة زاوية واحدة، كما يدور جسم صلب،
وتعبره النجوم والغاز؛ ولذلك لا يلتفّ شكل
الأذرع.
\end{solution}

\begin{solution}{exo:b3:milky-way-galaxies:12}
(أ)~$a = (4.2\times10^6\times16^2)^{1/3} = \qty{1024}{au}$؛ عند
\qty{8178}{pc} يقابلها زاوية $1024/8178 = \ang{;;0.125}$ (الوحدة الفلكية عند فرسخ فلكي واحد
تقابل ثانية قوسية واحدة). (ب)~$R_S = 2GM/c^2 = 2\times\num{6.674e-11}
\times\num{8.35e36}/(\num{2.998e8})^2 = \qty{1.24e10}{m} = \qty{0.083}{au}$؛
يقابل القطر $2R_S$ زاوية $0.166/8178 = \num{2.0e-5}$ ثانية قوسية، أي نحو
\qty{20}{\micro as}. (ج)~$R_S = \qty{4.02e-7}{pc}$، و$M/(\tfrac43\pi
R_S^3) = \num{1.5e25}$ كتلةً شمسية لكلّ فرسخ فلكي مكعّب، أربعين ضعفًا
من الحدّ الأدنى الذي تفرضه حركة الرامي~A* وحجمه: يقع الحدّ على بُعد
عامل أربعين من كثافة الثقب الأسود نفسه.
\end{solution}

\begin{solution}{pb:b3:milky-way-galaxies:1}

\textbf{الجزء الأول.}

\textbf{1.} $v_r = A\,d\sin2\ell$: $+A d$ عند \ang{45}، 0 عند \ang{90}، $-A d$
عند \ang{135}، 0 عند \ang{180}.

\textbf{2.} $A = (15.4 + 15.2)/(2\times\qty{1.00}{kpc}) =
\qty{15.3}{km.s^{-1}.kpc^{-1}}$.

\textbf{3.} عند \ang{90}، $v_t = d(B - A) = 4.74\times(-5.91)\times1.00 =
\qty{-28.02}{km.s^{-1}}$، ومنه $B = 15.3 - 28.02 =
\qty{-12.72}{km.s^{-1}.kpc^{-1}}$.

\textbf{4.} عند \ang{0}، $v_t = d(A + B) = 4.74\times0.55 =
\qty{2.61}{km.s^{-1}}$، ومنه $B = 2.61 - 15.3 = \num{-12.69}$: القيمتان
تتّفقان ضمن \num{0.03}.

\textbf{5.} $\omega_0 = A - B = \qty{28.0}{km.s^{-1}.kpc^{-1}}$
و$V_0 = 28.0\times8.18 = \qty{229}{km.s^{-1}}$.

\textbf{6.} $V'_0 = -(A + B) = \qty{-2.6}{km.s^{-1}.kpc^{-1}}$: يتناقص برفق.

\textbf{7.} $P = 2\pi/\omega_0$: $2\pi\times\qty{8.18}{kpc}/\qty{229.2}{km.s^{-1}}
= \qty{219}{Myr}$؛ $4540/219 = 20.7$ دورة.

\textbf{الجزء الثاني.}

\textbf{8.} تحتاج نقطة المماسّ إلى خط النظر الذي يمرّ أقرب إلى
المركز من الشمس، $\sin\ell > 0$ مع $\ell < \ang{90}$ (أو الربع
المتناظر، $\ang{270} < \ell < \ang{360}$، حيث تكون أكبر
سرعة هي الأكثر سلبيةً). خارج الدائرة الشمسية تتزايد المسافة
إلى المركز على امتداد خط النظر كلّه:
فلا تُفرَد نقطة بعينها، ولا تُخبر السرعة وحدها أين يقع الغاز.

\textbf{9.} $R = 8.18\sin\ell$: \qtylist{5.26;6.27;7.08;7.69}{kpc}.

\textbf{10.} $V = v_\mathrm{max} + V_0\sin\ell$: $82.9 + 147.2 = 230.1$,
$55.6 + 175.4 = 231.0$, $32.4 + 198.3 = 230.7$, $14.9 + 215.2 =
\qty{230.1}{km.s^{-1}}$.

\textbf{11.} $\Delta\nu = \nu v/c = \num{1.420406e9}\times82.9/299792 =
\qty{393}{kHz}$، نحو ترددات أدنى: الغاز يبتعد.

\textbf{12.} ضمن نحو \qty{3}{km.s^{-1}} من العمالقة الحمراء: بين 5
و\qty{8}{kpc} يكون \omterm{def:b3:milky-way-galaxies:rotation-curve}{منحنى الدوران} مستويًا.

\textbf{13.} تُعطي الكتلة النقطية $V = V_0\sqrt{R_0/R} = 229\times\sqrt{8.18/5.26} =
\qty{286}{km.s^{-1}}$ عند نقطة المماسّ الأولى، وهي أسرع بكثير من
\qty{230}{km.s^{-1}} المستخلصة: الكتلة غير متمركزة عند المركز،
ومعظمها يقع بين 5 و\qty{8}{kpc} وما وراءها.

\textbf{الجزء الثالث.}

\textbf{14.} $\num{1.067e10} + 2\times\num{3.944e10} = \num{8.96e10}$ كتلةً
شمسية، أي نحو ضعف تقدير المراجعة \num[separate-uncertainty=true]{5 \pm 1 e10}؛ ربما تكون نجوم النموذج أثقل من اللازم،
ولو صحّ ذلك لجعل المادة المظلمة المستخلصة أدناه
تقديرًا ناقصًا.

\textbf{15.} $V_b^2 = GMR^2/(R^2 + b^2)^{3/2} = \num{4.301e-6}\times\num{1.067e10}
\times400/(400.09)^{3/2}$, $V_b = \qty{47.9}{km.s^{-1}}$.

\textbf{16.} القرص الرقيق: $(a + b)^2 = 30.80$، $V^2 = \num{4.301e-6}\times\num{3.944e10}
\times400/(430.8)^{3/2}$، $V = \qty{87.1}{km.s^{-1}}$. \omterm{def:b3:milky-way-galaxies:structure}{القرص السميك}:
$(a + b)^2 = 11.56$, $V = \qty{90.1}{km.s^{-1}}$.

\textbf{17.} $V_\mathrm{vis} = \sqrt{47.9^2 + 87.1^2 + 90.1^2} =
\qty{134.2}{km.s^{-1}}$.

\textbf{18.} عند \qty{8.18}{kpc}: 74.8، 108.4 و\qty{127.8}{km.s^{-1}}،
مجتمعةً \qty{183.5}{km.s^{-1}}، مقابل $V_0 = \qty{229}{km.s^{-1}}$:
تقصر الكتلة المرئية عن المطلوب حتى عند الشمس.

\textbf{19.} $\sqrt{G\times\num{8.96e10}/20} = \qty{138.8}{km.s^{-1}}$.
قيمة النموذج \qty{134.2}{km.s^{-1}} قريبة منها، لأنّ عند \qty{20}{kpc}،
على بُعد أربعة إلى ستة أطوال مقياس القرص، تقع معظم الكتلة المرئية
داخل المدار: صار المنحنى المرئي كيبلريًّا.

\textbf{الجزء الرابع.}

\textbf{20.} $M = 20\times200^2/\num{4.301e-6} = \num{1.860e11}$ كتلةً شمسية.

\textbf{21.} $M_\mathrm{vis} = 20\times134.2^2/\num{4.301e-6} = \num{8.37e10}$
كتلةً شمسية.

\textbf{22.} $M_\mathrm{dark} = \num{1.860e11} - \num{8.37e10} = \num{1.02e11}$;
لأنّ الكتلتين كلتيهما تساويان $RV^2/G$ عند نفس $R$، فإنّ الكسر هو
$1 - (V_\mathrm{vis}/V)^2 = 1 - (134.2/200)^2 = \num{0.550}$.

\textbf{23.} $x = 20/14.8 = 1.351$, $\ln(2.351) - 1.351/2.351 = 0.8549 - 0.5748
= 0.2801$, $M = 4\pi\times\num{1.06e7}\times14.8^3\times0.2801 =
\num{1.21e11}$ كتلةً شمسية؛ ونسبة المادة المظلمة هي
$1.21/(1.21 + 0.837) = \num{0.59}$. تُطابَق \omterm{def:b3:milky-way-galaxies:halo}{هالة المادة المظلمة} على المنحنى كلّه،
الذي تتجاوزه الهالة عند \qty{20}{kpc}: تنخفض القياسات هناك
بنحو \qty{10}{km.s^{-1}} دون النموذج الأملس (وهو \qty{209}{km.s^{-1}}).

\textbf{24.} $x = 8.18/14.8 = 0.553$, $\rho = \num{1.06e7}/(0.553\times1.553^2)
= \num{7.95e6}$ كتلةً شمسية لكلّ كيلوفرسخ فلكي مكعّب، أي \num{0.0080} لكلّ
فرسخ فلكي مكعّب، أو \qty{0.30}{GeV.cm^{-3}}، وهي القيمة التي يذكرها مؤلّفو
المطابقة.

\textbf{25.} $1 - 8.37/18.60 = \num{0.550}$: نسبة المادة المظلمة من
الكتلة الواقعة ضمن \qty{20}{kpc} هي
\textbf{\qty[text-series-to-math=true]{55}{\percent}}.
\end{solution}
